\section{Capitoli 12-14: Liste Concatenate}

\begin{dizionario}{Lista vs Array}
\textbf{Lista:} Nodi sparsi collegati da punti `next`. Accesso $O(N)$.
\textbf{Array:} Blocco contiguo. Accesso $O(1)$.
\textbf{Nodo:} Oggetto `{ val: ..., next: ... }`.
\end{dizionario}

\begin{esercizio}{Es. 29: Conta Occorrenze (Lista)}
\textbf{TRACCIA:} Scrivere `countVal(nodo, target)` che conta quante volte `target` appare nella lista.
\end{esercizio}

\begin{jsbox}
function countVal(nodo, target) {
    if (!nodo) return 0;
    
    let match = (nodo.val === target) ? 1 : 0;
    return match + countVal(nodo.next, target);
}
\end{jsbox}

\begin{esercizio}{Es. 30: Ultimo Elemento (Il 'Last')}
\textbf{TRACCIA:} Restituire il valore dell'ultimo nodo della lista.
\textit{Attenzione: non l'intera lista, solo il valore.}
\end{esercizio}

\begin{jsbox}
function getLast(nodo) {
    if (!nodo) return null; // Lista vuota
    if (!nodo.next) return nodo.val; // Sono l'ultimo!
    return getLast(nodo.next);
}
\end{jsbox}

\begin{esercizio}{Es. 31: Filtra Lista (List Filter)}
\textbf{TRACCIA:} Scrivere `listFilter(nodo, predicato)` che restituisce una **nuova** lista contenente solo i nodi che passano il test.
\end{esercizio}

\begin{jsbox}
function listFilter(nodo, predicato) {
    if (!nodo) return null;
    
    // Se il nodo corrente passa, lo tengo e continuo
    if (predicato(nodo.val)) {
        return {
            val: nodo.val,
            next: listFilter(nodo.next, predicato)
        };
    } else {
        // Se non passa, lo salto (ritorno direttamente il risultato del prossimo)
        return listFilter(nodo.next, predicato);
    }
}
\end{jsbox}

\begin{esercizio}{Es. 32: Append (Concatena Due Liste)}
\textbf{TRACCIA:} Scrivere `listConcat(l1, l2)` che attacca `l2` in fondo a `l1`.
\end{esercizio}

\begin{jsbox}
function listConcat(l1, l2) {
    if (!l1) return l2; // Se l1 è finita, attacco l2
    
    return {
        val: l1.val,
        next: listConcat(l1.next, l2)
    };
}
\end{jsbox}

\begin{esercizio}{Es. 42: Inverti Lista (Reverse)}
\textbf{TRACCIA:} Scrivere `reverseList(head)` che inverte i puntatori della lista e restituisce la nuova testa.
\textit{Iterativo o Ricorsivo.}
\end{esercizio}

\begin{jsbox}
function reverseList(head) {
    let prev = null;
    let curr = head;
    
    while (curr !== null) {
        let nextTemp = curr.next; // Salvo il prossimo
        curr.next = prev;         // Inverto il puntatore
        prev = curr;              // Avanzo prev
        curr = nextTemp;          // Avanzo curr
    }
    return prev; // Nuova testa
}
\end{jsbox}

\begin{esercizio}{Es. 43: Rileva Ciclo (Floyd)}
\textbf{TRACCIA:} Determinare se una lista ha un ciclo (un nodo punta a uno precedente).
\textit{Usa due puntatori: 'tartaruga' (lento) e 'lepre' (veloce).}
\end{esercizio}

\begin{jsbox}
function hasCycle(head) {
    let slow = head;
    let fast = head;
    
    while (fast && fast.next) {
        slow = slow.next;       // 1 passo
        fast = fast.next.next;  // 2 passi
        
        if (slow === fast) return true; // Si sono incontrati!
    }
    return false;
}
\end{jsbox}
